\chapter{Installation and launch}
\label{chap:installation}

\section{Prerequisites}

\begin{itemize}
    \item \textbf{Operating system} : Windows 10 or 11
        (tested), Linux and macOS (compatibility expected, not
        tested by the developers).
    \item \textbf{Python} : version 3.10 or higher for the XFoil
        solver (recommendation: 3.13). The optional \emph{FlexFoil}
        solver specifically requires \textbf{Python 3.11} (only
        \code{cp311} wheels are published).
    \item \textbf{XFoil solver} : required only for
        running simulations. On Windows, the executable
        \chemin{xfoil.exe} is provided in the
        \chemin{externaltools/xfoil/} folder of the distribution.
\end{itemize}

\section{Windows executable distribution}

For Windows users who do not wish to install Python,
a standalone \chemin{.exe} distribution is available. It
bundles the Python interpreter, all dependencies and the
XFoil binary.

The application is launched by double-clicking on
\chemin{AirfoilTools.exe}.

\begin{noteinfo}
On first launch, Windows may display a
\emph{SmartScreen} warning because the executable is not signed. Click
on \og More info \fg{} then \og Run
anyway \fg{} to allow it to launch.
\end{noteinfo}

\section{Installation from source}

\subsection{Retrieving the repository}

\begin{lstlisting}[language=bash,caption={Cloning the repository}]
git clone https://github.com/BenoitGFree/AirfoilTools.git
cd AirfoilTools
\end{lstlisting}

\subsection{Creating the Python environment}

\begin{lstlisting}[language=bash,caption={Creating the venv and installing
the dependencies (Windows)}]
py -3 -m venv env_py3
env_py3\Scripts\activate
pip install -r requirements.txt
\end{lstlisting}

The \chemin{requirements.txt} file lists the dependencies :

\begin{itemize}
    \item \code{numpy} --- vector computation ;
    \item \code{scipy} --- numerical algorithms (interpolation,
        optimization) ;
    \item \code{matplotlib} --- 2D plotting ;
    \item \code{PySide6} --- Qt for Python graphical interface.
\end{itemize}

\subsection{Launching the application}

\begin{lstlisting}[language=bash,caption={Starting the GUI}]
env_py3\Scripts\python.exe run_gui.py
\end{lstlisting}

The application opens on the main window shown in
\figref{01_main_window.png}.

\section{Verifying the installation}

To verify that the environment is correctly configured, it
is possible to run the unit test suite :

\begin{lstlisting}[language=bash,caption={Running the tests}]
env_py3\Scripts\python.exe -c "import unittest; import sys; ^
sys.path.insert(0, 'sources'); sys.path.insert(0, 'tests'); ^
from test_bezier import *; ^
unittest.main(module='test_bezier', exit=False, verbosity=2)"
\end{lstlisting}

The full suite contains more than 240 tests, all of which must
pass successfully.

\begin{attention}
Automatic detection of the XFoil executable depends on
the location of \chemin{xfoil.exe}. The application looks for it
in this order :
\begin{enumerate}
    \item next to the application executable (PyInstaller
        \emph{frozen} mode) ;
    \item in the \chemin{externaltools/xfoil/} folder ;
    \item in the system PATH.
\end{enumerate}
If XFoil is not found, the \menu{XFoil Settings} tab
works but launching a simulation will fail.
\end{attention}
